%% ============================================================
%%  Section 5 — Computational Experiments and Results
%%  Article: BRKGA with Rolling Horizon for the MDRP
%%  Citation style: LNCS (numeric \cite{})
%% ============================================================

\section{Computational Experiments and Results}\label{sec:results}

%% ------------------------------------------------------------------
\subsection{Experimental Setup}\label{subsec:setup}
%% ------------------------------------------------------------------

\paragraph{Instances.}
We evaluate the four methods on 24 real-world instances
drawn from a Latin American food-delivery platform.
Each instance corresponds to a full operating day and
contains orders, couriers, and restaurant locations
from one or more urban zones.
Following the classification used in related work, instances are
grouped by total order volume into three tiers:
\emph{Small} ($|O| < 5{,}000$, 13 instances),
\emph{Medium} ($5{,}000 \leq |O| \leq 20{,}000$, 5 instances), and
\emph{Large} ($|O| > 20{,}000$, 6 instances).
Instance sizes range from 538 to 59{,}167 orders and from
396 to 16{,}782 active couriers.

\paragraph{Methods compared.}
We compare four algorithms under identical rolling-horizon
conditions ($\Delta = 2$ min, horizon 0--1\,440 min):
\begin{enumerate}
    \item \textbf{BRKGA-Coverage} — BRKGA-RH with coverage-maximising
          fitness $f_{\text{cov}}$ (Section~\ref{subsec:brkga_cov}).
    \item \textbf{BRKGA-Travel} — BRKGA-RH with travel-minimising
          fitness $f_{\text{trav}}$ (Section~\ref{subsec:brkga_trav}).
    \item \textbf{Business Model} — deterministic rolling-horizon
          greedy without lookahead (Section~\ref{subsubsec:maxcover_base}).
    \item \textbf{Business Model Lookahead} — Business Model extended
          with demand-forecast deferral (Section~\ref{subsubsec:maxcover_la}).
\end{enumerate}

\paragraph{Evaluation metrics.}
We report two primary metrics for solution quality:
(i) the \emph{coverage rate} $\rho = |\bar{O}|/|O|$ (higher is better),
where $\bar{O}$ is the set of orders successfully delivered by the
\emph{end of the full planning day} (not per window), and
(ii) the \emph{average travel time per delivered order} $\bar{\tau}$
in minutes (lower is better).
Note that inside the BRKGA decoder a \emph{per-window} coverage
$\rho^w = |\mathcal{A}^w|/|O^w|$ is used as part of the fitness
functions $f_{\text{cov}}$ and $f_{\text{trav}}$ to guide search
within each 2-minute window; the end-of-day $\rho$ accumulates these
decisions over the full 1,440-minute horizon.
We also report \emph{total daily execution time} (sum over all
windows) in seconds as the primary efficiency metric;
see Table~\ref{tab:results_all} caption for per-window context.
All experiments were run on commodity hardware
(Intel Core i7, 16 GB RAM, Windows 11).

%% ------------------------------------------------------------------
\subsection{Results by Instance}\label{subsec:results_table}
%% ------------------------------------------------------------------

Table~\ref{tab:results_all} reports, for each instance, the coverage
rate $\rho$ (\%) and average travel time $\bar{\tau}$ (min/order)
for BRKGA-Coverage and BRKGA-Travel, the coverage rate for both
Business Model variants, and the total execution time (s) for all
methods except the plain Business Model (which completes in under
60 s on every instance).
Size codes: Small (S), Medium (M), Large (L).
Bold marks the best $\rho$ per instance.

\begin{table}[ht]
\centering
\caption{Per-instance results for all four methods.
$\rho$ = coverage rate (\%) computed over the full planning day,
$\bar{\tau}$ = avg travel time (min/order),
$T$ = \textbf{total daily} execution time (s), i.e., summed over all
rolling-horizon windows of the 24-hour day.
The per-window budget is $\Delta = 120$\,s; daily totals can exceed
this because up to 720 windows are solved per instance.
Best $\rho$ per row in bold.}
\label{tab:results_all}
\scriptsize
\setlength{\tabcolsep}{2.8pt}
\begin{tabular}{rrlr|cc|cc|c|c|rrr}
\toprule
 & & & &
\multicolumn{2}{c|}{\textbf{BRKGA-Cov}} &
\multicolumn{2}{c|}{\textbf{BRKGA-Trav}} &
\textbf{BM} &
\textbf{BML} &
\multicolumn{3}{c}{\textbf{Exec.\ time (s)}} \\
\textbf{Inst} & \textbf{Sz} & $|O|$ & $|C|$ &
$\rho$ & $\bar{\tau}$ & $\rho$ & $\bar{\tau}$ &
$\rho$ & $\rho$ &
$T_{\text{BC}}$ & $T_{\text{BT}}$ & $T_{\text{BML}}$ \\
\midrule
0  & L & 56{,}420 & 16{,}710 & 37.1 & 34.5 & 37.4 & 34.5 & 41.0 & \textbf{46.3} & 1{,}143 &    571 & 2{,}916 \\
1  & S &  2{,}959 &  1{,}185 & \textbf{57.7} & 22.8 & \textbf{57.7} & 22.8 & 33.7 & 40.4 &      20 &     19 &    38 \\
2  & L & 24{,}720 &  6{,}570 & \textbf{59.0} & 29.4 & 59.0 & 29.4 & 38.5 & 44.2 &     458 &    356 &   458 \\
3  & S &  2{,}959 &  1{,}185 & \textbf{57.3} & 22.6 & 57.1 & 22.5 & 33.7 & 40.4 &      23 &     23 &    37 \\
4  & S &    844   &    415   & \textbf{54.1} & 24.7 & \textbf{54.1} & 24.8 & 37.7 & 44.1 &       6 &      6 &    11 \\
5  & M &  9{,}115 &  2{,}803 & 48.8 & 22.2 & \textbf{48.9} & 22.2 & 26.0 & 33.0 &      80 &     80 &    98 \\
6  & L & 59{,}167 & 16{,}782 & 34.4 & 34.4 & 34.0 & 34.5 & 42.8 & \textbf{48.7} &     490 &    477 & 6{,}738 \\
7  & M &  7{,}357 &  2{,}643 & \textbf{58.2} & 31.1 & 57.9 & 31.1 & 36.1 & 44.9 &     116 &     12 &   103 \\
8  & L & 23{,}155 &  6{,}851 & \textbf{56.3} & 31.3 & 56.2 & 31.3 & 38.2 & 46.0 &      79 &     82 &   436 \\
9  & S &  2{,}407 &  1{,}156 & \textbf{60.7} & 24.2 & \textbf{60.7} & 24.0 & 39.1 & 46.4 &       2 &      3 &    29 \\
10 & S &    735   &    455   & \textbf{62.0} & 23.7 & \textbf{62.0} & 23.8 & 46.3 & 53.6 &       5 &      5 &     8 \\
11 & M &  8{,}451 &  3{,}023 & 50.4 & 27.4 & \textbf{50.6} & 27.8 & 28.1 & 34.6 &     110 &     16 &   109 \\
12 & L & 31{,}199 & 13{,}152 & 51.6 & 34.3 & 51.4 & 34.0 & 51.9 & \textbf{59.2} &     199 &    304 & 1{,}201 \\
13 & S &  4{,}516 &  1{,}876 & 70.1 & 30.5 & \textbf{70.1} & 30.8 & 51.2 & 62.0 &      56 &      6 &    79 \\
14 & M & 13{,}051 &  4{,}824 & 68.2 & 32.1 & \textbf{68.4} & 32.3 & 50.0 & 57.9 &      32 &     36 &   360 \\
15 & S &  1{,}462 &    941   & \textbf{66.4} & 23.2 & \textbf{66.4} & 23.1 & 48.2 & 57.1 &      17 &      2 &    32 \\
16 & S &    538   &    430   & \textbf{70.8} & 23.6 & \textbf{70.8} & 23.7 & 63.0 & 66.7 &       4 &      4 &     8 \\
17 & S &  4{,}390 &  2{,}368 & \textbf{66.4} & 27.1 & 66.4 & 27.0 & 44.3 & 51.4 &       6 &      6 &    73 \\
18 & L & 30{,}448 & 12{,}218 & 45.1 & 33.8 & 46.2 & 34.0 & 56.0 & \textbf{62.4} &     295 &    459 &   966 \\
19 & S &  4{,}389 &  1{,}908 & \textbf{72.0} & 30.6 & 72.0 & 30.8 & 53.5 & 63.5 &      57 &      6 &    69 \\
20 & M & 13{,}076 &  4{,}905 & 69.7 & 31.1 & \textbf{69.8} & 31.2 & 51.5 & 59.5 &      31 &     34 &   266 \\
21 & S &  1{,}376 &    959   & \textbf{64.8} & 23.0 & \textbf{64.8} & 22.8 & 43.9 & 52.8 &      15 &      1 &    22 \\
22 & S &    539   &    396   & \textbf{74.0} & 23.6 & \textbf{74.0} & 23.7 & 61.0 & 68.7 &       3 &      3 &     6 \\
23 & S &  4{,}387 &  2{,}111 & 65.7 & 27.5 & \textbf{65.9} & 27.5 & 42.5 & 50.3 &      53 &      4 &    69 \\
\midrule
\multicolumn{4}{r|}{\textbf{Average}} &
59.2 & 27.9 & \textbf{59.2} & 27.9 & 44.1 & 51.4 & 138 & 105 & 589 \\
\bottomrule
\end{tabular}
\end{table}

%% ------------------------------------------------------------------
\subsection{Coverage Rate Analysis}\label{subsec:coverage_analysis}
%% ------------------------------------------------------------------

\paragraph{Overall rankings.}
Averaged over all 24 instances, BRKGA-Coverage and BRKGA-Travel
achieve virtually identical coverage rates of \textbf{59.2\%},
followed by Business Model Lookahead at \textbf{51.4\%} and
Business Model at \textbf{44.1\%}.
The two BRKGA variants thus outperform the Business Model Lookahead
by 7.8 percentage points on average, and the Business Model by
15.1 percentage points.

\paragraph{Effect of instance size.}
The relative performance of the methods varies substantially with
instance size, as summarised in Table~\ref{tab:results_by_size}.

\begin{table}[ht]
\centering
\caption{Average coverage rate $\rho$ (\%) by instance size tier.}
\label{tab:results_by_size}
\small
\begin{tabular}{lrrrrrr}
\toprule
\textbf{Size} & \textbf{Count} & \textbf{Avg $|O|$} &
\textbf{BRKGA-Cov} & \textbf{BRKGA-Trav} &
\textbf{BML} & \textbf{BM} \\
\midrule
Small  & 13 & 2{,}423 & 64.8\% & 64.8\% & 53.6\% & 46.0\% \\
Medium &  5 & 10{,}210 & 59.1\% & 59.1\% & 46.0\% & 38.4\% \\
Large  &  6 & 37{,}518 & 47.2\% & 47.3\% & 51.1\% & 44.8\% \\
\midrule
\textbf{All} & \textbf{24} & \textbf{11{,}688} & \textbf{59.2\%} & \textbf{59.2\%} & \textbf{51.4\%} & \textbf{44.1\%} \\
\bottomrule
\end{tabular}
\end{table}

On \emph{Small} and \emph{Medium} instances the BRKGA variants
maintain a consistent advantage of 11--20 percentage points over
the Business Model Lookahead, which in turn outperforms the
Business Model by 7--8 percentage points.
However, the picture reverses on \emph{Large} instances:
the Business Model Lookahead achieves an average of \textbf{51.1\%},
surpassing BRKGA-Coverage (47.2\%) by 3.9 percentage points
and the Business Model (44.8\%) by 6.3 percentage points.
This reversal occurs on all six large instances and is most
pronounced on Instances 0 and 6
(56{,}420 and 59{,}167 orders, respectively),
where BRKGA-Coverage drops to 37.1\% and 34.4\% while
Business Model Lookahead reaches 46.3\% and 48.7\%.

The explanation lies in scalability:
on large instances the BRKGA's $G_{\max} = 10$ generations and
$N_p = 10$ population are insufficient to explore the solution
space effectively within the 2-minute window budget,
whereas the Business Model Lookahead's deterministic two-phase
approach scales gracefully and benefits from the deferral
mechanism to consolidate larger, higher-quality batches.

\paragraph{Lookahead benefit.}
Across all 24 instances the Business Model Lookahead improves
upon the Business Model by an average of \textbf{7.3 percentage
points} (range: 3.7--10.8 pp), confirming that demand-forecast
deferral consistently yields better batch utilisation at the
cost of slightly delayed dispatch.

%% ------------------------------------------------------------------
\subsection{Travel Time Analysis}\label{subsec:travel_analysis}
%% ------------------------------------------------------------------

The average travel time per delivered order is remarkably similar
between BRKGA-Coverage and BRKGA-Travel: \textbf{27.86 vs
27.90 min/order} overall, with differences below 0.5 min/order
on every instance.
This near-parity reveals an important structural property
of the problem: the nearest-neighbour decoder produces routes
whose travel efficiency is largely determined by geographic
density and time-window constraints rather than by the
fitness-function penalty weights.
In other words, once a feasible batch is assembled, its
travel time is almost entirely dictated by the delivery
sequence heuristic, leaving little room for the fitness function
to differentiate between the two BRKGA variants.

Travel times are systematically shorter on Small instances
($\bar{\tau} \approx 25.0$ min/order) than on Large ones
($\bar{\tau} \approx 32.9$ min/order), reflecting the wider
geographic spread of deliveries in high-volume urban zones.

%% ------------------------------------------------------------------
\subsection{Computational Efficiency}\label{subsec:time_analysis}
%% ------------------------------------------------------------------

Table~\ref{tab:results_all} shows substantial differences in
total execution time across the methods.
BRKGA-Coverage averages \textbf{138 s} per instance
(range: 2--1{,}143 s) and BRKGA-Travel \textbf{105 s}
(range: 1--571 s), while Business Model Lookahead requires
on average \textbf{589 s} (range: 6--6{,}738 s).
The plain Business Model completes in under 60 s on every instance.

Contrary to what might be expected, the two BRKGA variants are
faster in total execution time than the Business Model Lookahead.
BRKGA-Travel is the fastest method overall, running
\textbf{4.3$\times$ faster} than BML on average (105 s vs 589 s),
and BRKGA-Coverage is \textbf{4.3$\times$ faster} as well (138 s vs 589 s).
This occurs because each BRKGA window runs at most $G_{\max}=10$
generations over a population of $N_p=10$, producing a compact
per-window footprint, whereas the Business Model Lookahead's
forecast computations and simulation-level bookkeeping accumulate
over the full operating day.
BRKGA-Travel is also on average \textbf{24\% faster} than
BRKGA-Coverage (105 s vs 138 s), suggesting that the
travel-minimising fitness induces a slightly simpler search
landscape without sacrificing coverage quality.

The BML overhead grows sharply with instance size:
on Instance~6 (59{,}167 orders) it reaches 6{,}738 s (1.87 h),
versus 490 s for BRKGA-Coverage and 477 s for BRKGA-Travel.
On small instances the difference is also notable
(e.g., Instance~9: BML 29 s vs BC 2 s).
All methods respect the per-window $\Delta = 2$-min dispatch
budget, so real-time delivery dispatch is guaranteed regardless
of total run time.

%% ------------------------------------------------------------------
\subsection{Discussion}\label{subsec:discussion}
%% ------------------------------------------------------------------

The results reveal a clear trade-off between solution quality and
computational cost that is strongly mediated by instance scale.

For \emph{Small and Medium} instances ($|O| \leq 20{,}000$),
the BRKGA variants are the preferred choice when execution time
is not a binding constraint: they consistently outperform both
Business Model variants by 11--20 percentage points on coverage
while producing comparable travel times.
The two BRKGA fitness functions yield nearly identical outcomes,
suggesting that the coverage-maximising objective ($f_{\text{cov}}$)
and the travel-minimising objective ($f_{\text{trav}}$) converge
to similar solutions under the tight 10-generation window budget.

For \emph{Large} instances ($|O| > 20{,}000$), the Business Model
Lookahead is the superior method in terms of \emph{coverage quality}:
it achieves higher coverage than both BRKGA variants on all six
large instances, and consistently satisfies the 2-minute per-window
budget.
However, the BRKGA variants remain faster in total run time even
on large instances (e.g., Instance~6: 490 s for BC vs 6{,}738 s for BML),
so the choice between methods involves a trade-off between
coverage quality and computational overhead.
The demand-forecast component is the key differentiator over the
plain Business Model, providing a 6--11 pp coverage improvement
by enabling larger, more efficient batches.

These findings suggest a \emph{hybrid deployment strategy}:
use BRKGA-Coverage for small-to-medium instances where the
evolutionary search has sufficient time to explore the solution
space, and switch to Business Model Lookahead for large instances
where the combinatorial complexity exceeds the BRKGA's effective
search budget.
